\section{AI SRE: el nuevo paradigma de las operaciones inteligentes}

\subsection{Herramientas principales}

\begin{itemize}
    \item \textbf{Resolve AI}: el invitado de esta semana, Mayank Agarwal, es su CTO
    \item \textbf{DataDog Bits AI Agent}: análisis de IA basado en los datos de monitoreo de DataDog
    \item \textbf{Splunk Observability Assistant}: el asistente de IA de observabilidad de Splunk
\end{itemize}

\subsection{Características centrales de AI SRE}

\begin{enumerate}
    \item \textbf{Grafo de conocimiento dinámico}: mapea de manera dinámica las dependencias y el estado entre los servicios
    \item \textbf{Sistema Agentic que abarca toda la pila de observabilidad}: opera en múltiples nubes y herramientas de monitoreo
    \item \textbf{Generación de narrativas en tiempo real}: genera una narrativa en tiempo real de lo que está ocurriendo, ubica posibles causas raíz con evidencia de respaldo y recomienda pasos concretos de corrección
    \item \textbf{Interpretabilidad y auditabilidad}: enfatiza la transparencia de las predicciones y el razonamiento
\end{enumerate}

\begin{knowledgebox}{AI SRE resuelve el problema de los silos de información}
Uno de los valores más importantes de AI SRE es romper los silos de conocimiento dentro de la organización. En el modelo tradicional, el conocimiento sobre dependencias no documentadas, servicios heredados frágiles y peculiaridades del sistema que solo se manifiestan durante incidentes de alta presión suele existir únicamente en la mente de unos cuantos ingenieros. AI SRE sistematiza ese conocimiento organizacional y a nivel de servicio, disperso, y lo pone a disposición de todo el equipo.
\end{knowledgebox}

\subsection{Cambios que trae la IA}

\begin{itemize}
    \item La automatización reduce el tiempo de revisión y detecta problemas de manera temprana
    \item Los desarrolladores aprenden las mejores prácticas a partir de las sugerencias de la IA
    \item La IA aplica los mismos estándares a todas las revisiones de código
    \item La IA se encarga de las verificaciones rutinarias, lo que permite que los humanos se concentren en la lógica compleja
    \item Los sistemas de IA mejoran de manera continua a medida que se acumulan datos
    \item Las herramientas modernas de revisión de código con IA ofrecen análisis contextual y reconocimiento de patrones
\end{itemize}

\subsection{Limitaciones actuales}

\begin{warningbox}{Limitaciones centrales de AI SRE}
\begin{itemize}
    \item La \textbf{complejidad} de los incidentes que se pueden resolver es limitada — por ahora la IA es más adecuada para problemas de patrones conocidos
    \item La \textbf{heterogeneidad} de las pilas tecnológicas de producción modernas dificulta la comprensión por parte de la IA
    \item La capacidad de \textbf{corregir el código real} a partir de los resultados de detección todavía está en una etapa temprana — por ahora todos los proveedores parten del análisis de causa raíz
    \item Un buen análisis de causa raíz depende de \textbf{una buena infraestructura de monitoreo} (jardinería de monitoreo/monitoring gardening)
    \item La seguridad puede convertirse en un nuevo vector de ataque
\end{itemize}
\end{warningbox}

\subsection{La evolución de sistemas de alertas a sistemas de operación}

Una de las percepciones clave de la semana 9 es que AI SRE no debe ser solamente un «resumidor de alertas más inteligente», sino que debe evolucionar de manera gradual hacia un verdadero sistema de colaboración operativa.

\begin{itemize}
    \item La primera etapa es la \textbf{agregación de información}: reunir bitácoras, métricas, traces y registros de cambios en una sola vista
    \item La segunda etapa es la \textbf{inferencia de causa raíz}: presentar cadenas causales y evidencia que se asemejen más al razonamiento de un ingeniero
    \item La tercera etapa es la \textbf{recomendación de ejecución}: sugerir acciones de mitigación, rutas de reversión y plantillas de comunicación
    \item A más largo plazo llega la \textbf{corrección automática controlada}: ejecutar de manera automática acciones de corrección de bajo riesgo dentro de límites claramente definidos
\end{itemize}

\begin{knowledgebox}{Un sistema AI SRE maduro primero debe aprender a ``cometer menos errores''}
En un entorno de producción, una automatización equivocada suele ser más peligrosa que una operación humana un poco más lenta. Por eso la madurez de AI SRE no se mide solo por cuántas cosas puede automatizar, sino por si puede identificar de manera constante cuáles no debe automatizar.
\end{knowledgebox}

\subsection{Resumen de la sección}

AI SRE apoya las operaciones mediante grafos de conocimiento dinámicos, sistemas Agent que abarcan toda la pila y generación de narrativas en tiempo real. Su mayor valor consiste en romper los silos de información y reducir el MTTR. Su limitación actual es que solo puede manejar problemas de patrones conocidos con complejidad media, y todavía queda un largo camino de la detección a la corrección automática.

